%! AP Statistics — Unit 5, Lesson 5.2 — Guided Notes
\documentclass[10pt]{article}
\usepackage{apstats-article}
\usepackage{apstats-boxes}
\newcommand{\TallMath}[1]{$\displaystyle #1\rule[-1.4em]{0pt}{3.2em}$}

\begin{document}

\pageheader{Unit 5, Lesson 5.2}{Guided Notes}
\namedateperiod

\vspace{0.15in}

\begin{objectivebox}
By the end of this lesson, I will be able to:
\begin{enumerate}
  \item \writeline
  \item \writeline
  \item \writeline
\end{enumerate}
\end{objectivebox}

\vspace{0.1in}

\begin{vocabbox}
\termblanklong{Continuous random variable}
\termblanklong{Normal distribution}
\termblanklong{z-score}
\termblanklong{Standard normal distribution}
\end{vocabbox}

\vspace{0.1in}

\begin{hookbox}
Adult male heights $\sim N(69.2, 3.6)$. Estimate the fraction of adult males taller
than 6 feet (72 in) from the histogram (shown on board), then we will calculate exactly.

My estimate: \blank{3cm} \qquad Calculated answer: \blank{3cm}
\end{hookbox}

\vspace{0.1in}

\begin{notesbox}{Normal Distributions: Key Facts}
Notation: $X \sim N(\mu, \sigma)$ means $X$ is normally distributed with mean
\blank{3cm} and standard deviation \blank{3cm}.

The \textbf{area} under the curve over an interval equals \blank{6cm} (VAR-6.A.3).

\textbf{z-score formula:} $z = $ \blank{4cm} \hfill
means $z$ is the number of \blank{4cm} that $x$ is from the mean.

\textbf{68--95--99.7 Rule:}
\begin{itemize}
  \item $P(\mu - \sigma < X < \mu + \sigma) \approx $ \blank{2cm}
  \item $P(\mu - 2\sigma < X < \mu + 2\sigma) \approx $ \blank{2cm}
  \item $P(\mu - 3\sigma < X < \mu + 3\sigma) \approx $ \blank{2cm}
\end{itemize}
\end{notesbox}

\vspace{0.1in}

\begin{notesbox}{Calculating Normal Probabilities}
\textbf{Forward probability} (area given a boundary):
\begin{enumerate}
  \item Sketch and shade the curve. Label \blank{3cm} and \blank{3cm}.
  \item Compute the z-score: $z = \frac{x - \mu}{\sigma} = $ \blank{4cm}
  \item Use table or \texttt{normalcdf(lower, upper, $\mu$, $\sigma$)} $=$ \blank{3cm}
  \item Interpret in context: ``The probability that \ldots'' \writeline
\end{enumerate}

\vspace{0.06in}
\textbf{Example:} $X \sim N(69.2, 3.6)$. Find $P(X > 72)$.\\[4pt]
$z = \frac{72 - 69.2}{3.6} = $ \blank{2.5cm} \qquad
$P(X > 72) = 1 - P(Z < \blank{1.5cm}) = 1 - $ \blank{2cm} $\approx$ \blank{2cm}
\end{notesbox}

\vspace{0.1in}

\begin{notesbox}{Finding a Boundary Value (Inverse Normal)}
\textbf{Inverse normal} (boundary given an area):
\begin{enumerate}
  \item Sketch and shade the curve. Label the given \blank{3cm}.
  \item Find the z-score: $z^* = $ \blank{3cm} (from table or \texttt{invNorm}).
  \item Solve for $x$: $x = \mu + z^*\sigma = $ \blank{6cm}
\end{enumerate}

\vspace{0.06in}
\textbf{Example:} $X \sim N(69.2, 3.6)$. Find the 90th percentile.\\[4pt]
$z^* = $ \blank{2.5cm} \qquad
$x = 69.2 + (\blank{2cm})(3.6) = $ \blank{3cm} inches
\end{notesbox}

\vspace{0.1in}

\begin{notesbox}{When Is the Normal Model Appropriate? (VAR-6.C)}
The normal model is appropriate when the distribution is \blank{6cm}
and \blank{8cm}.

Circle the correct answer for each:
\begin{itemize}
  \item Heights of adult males \hfill (appropriate / not appropriate)
  \item Number of heads in 5 coin flips \hfill (appropriate / not appropriate)
  \item Average of 100 random observations (by CLT) \hfill (appropriate / not appropriate)
\end{itemize}
\end{notesbox}

\vspace{0.1in}

\begin{practicebox}
$X \sim N(500, 100)$ (SAT Math scores).
\begin{enumerate}
  \item Find $P(X > 650)$. Show all steps including a sketch.

        \vspace{1.0in}

  \item Find the SAT Math score at the 25th percentile.

        \vspace{0.8in}
\end{enumerate}
\end{practicebox}

\end{document}
